\documentclass[11pt,letterpaper]{article}
\usepackage[T1]{fontenc}
\usepackage{lmodern,amsmath,amssymb,amsthm,microtype}
\usepackage[margin=0.9in]{geometry}
\usepackage[hidelinks]{hyperref}
\newtheorem{theorem}{Theorem}
\newtheorem{lemma}{Lemma}
\title{Rounding separated subset sums at the hypercube-bandwidth threshold}
\author{Aafi Khan Malek\\ \small Independent researcher, Ahmedabad, India\\
\small\href{mailto:aafimalek2032@gmail.com}{aafimalek2032@gmail.com}}
\date{28 November 2026}
\begin{document}
\maketitle
\begin{abstract}
Let \(H_n=\sum_{j=0}^{n-1}\binom{j}{\lfloor j/2\rfloor}\), the
bandwidth of the \(n\)-dimensional cube. For every positive integer \(n\), any
\(n\) positive real weights with unit-separated subset sums and
maximum below \(H_n+1\) round down to a set of positive integers
with distinct subset sums and maximum \(H_n\). Thus improving the
bandwidth lower bound by one is equivalent in the real and
integer problems. In particular the ten-weight real lower bound
is \(275\) using the reported integer computation. The exact real
optima are not determined.
\end{abstract}

\begin{theorem}\label{thm:round}
Suppose \(n\ge1\), and \(x_1<\cdots<x_n\) are positive reals with
pairwise unit-separated subset sums, and \(x_n<H_n+1\).
Then \(\lfloor x_1\rfloor,\ldots,\lfloor x_n\rfloor\) have
distinct subset sums and maximum \(H_n\).
Consequently, if \(f_{\mathbb R}(n)\) is the infimum of the real
maximum and \(f_{\mathbb Z}(n)\) is the minimum integer maximum,
\[
 f_{\mathbb R}(n)<H_n+1
 \iff f_{\mathbb R}(n)=H_n
 \iff f_{\mathbb Z}(n)=H_n .
\]
\end{theorem}

\section{An explicit saturated cut}
For \(n=1,2\), positive weights with unit-separated subset sums
give positive, distinct floors and \(x_n\ge n=H_n\), so the
theorem is immediate. Assume \(n\ge3\) henceforth.

Order the subsets of \([n]\) by cardinality and then lexicographically
by their increasing lists. This is the simplicial order in Harper's
vertex-isoperimetric theorem \cite{harper}. Write \(N=2^n\),
\(m=\lfloor n/2\rfloor\), \(r=\lfloor(n-1)/2\rfloor\), and set
\[
 A=\begin{cases}
 \{2,4,\ldots,n-2\},&n=2m,\\
 \{1,3,\ldots,n-2\},&n=2m+1.
 \end{cases}
\]
Let \(k\) be the one-based rank of \(A\). The initial segment through
\(A\) contains all smaller layers and a nonempty prefix of the
\(r\)-th layer. Its external boundary contains the rest of that layer
and precisely the prefix of the \((r+1)\)-st layer ending at
\(A\cup\{n\}\). To see the latter, the earliest \(r\)-subset of an
\((r+1)\)-set is obtained by deleting the largest element.
Its prefix must be at most \(A\), and the last \((r+1)\)-set
satisfying this is \(A\cup\{n\}\). The boundary therefore has size
\(\operatorname{rank}(A\cup\{n\})-k\).

For the arithmetic, write
\[
 W_m=\sum_{j=1}^{m}\binom{2j-1}{j-1},\qquad
 P_m=1+\sum_{j=1}^{m-1}\binom{2j+1}{j-1},\qquad
 S_m=1+\sum_{j=1}^{m-1}
 \left(\binom{2j+1}{j}-\binom{2j+1}{j-1}\right).
\]
The lexicographic rank of the alternating \(A\) within its layer
is \(P_m\): each skipped value one below a chosen value allows
\(j-1\) remaining choices out of \(2j+1\), for \(j=1,\dots,m-1\).
In the rank difference to \(A\cup\{n\}\), each skipped value
contributes the difference of binomial coefficients in \(S_m\).
The additional skipped last value \(n-1\) contributes one.
Including the one full layer in the difference gives
\[
 \operatorname{rank}(A\cup\{n\})-k=\binom nr+S_m=H_n.
\]
For the last identity, the equality at \(n=2\) is immediate.
Passing from \(2m\) to \(2m+1\) adds \(\binom{2m}{m}\) on
both sides by Pascal's rule. Passing from \(2m+1\) to \(2m+2\)
adds \(\binom{2m+1}{m}\), using the last summand
in \(S_{m+1}-S_m\).

The smaller full layers give
\[
 k=\begin{cases}
 N/2-\tfrac12\binom{2m}{m}-\binom{2m}{m-1}+P_m,& n=2m,\\
 N/2-\binom{2m+1}{m}+P_m,& n=2m+1.
 \end{cases}
\]
Since \(P_m+S_m=1+W_m\) and
\(\binom{2m}{m}/2=\binom{2m-1}{m-1}\), the integer
\(q=k+H_n-N/2\) equals \(1+W_{m-1}\) in even dimension and
\(1+W_m\) in odd dimension. Grouping consecutive central binomial
terms gives \(H_n-2q=W_m-1\) in both cases. Thus both
\([k,k+H_n]\) and its reflection by \(i\mapsto N+1-i\)
contain the middle ranks, and their union is \([k,N+1-k]\)
of rank span
\[
 L=N+1-2k=H_n+W_m.
\]
For \(n\ge4\), this is at least \(H_n+n-1\): \(W_2=4\) and
each further summand increases \(W_m\) by at least two.

\section{Rigid rounding inside the saturated slab}
Use ranks in the different order given by the actual weighted sums.
Write \(H=H_n\) and \(M=x_n\). A cube edge in coordinate \(i\)
has rank length at most \(\lfloor x_i\rfloor\). This is at most
\(H-1\) for \(i<n\), since singleton separation gives
\(x_i\le M-1<H\). By Harper's theorem and the explicit cut,
the first \(k\) ranks have at least \(H\) external neighbors.
All must lie in \([k+1,k+H]\). Rank \(k+H\) can only be
adjacent to rank \(k\) and only along coordinate \(n\). This
produces a largest-coordinate edge of real width \(M\) and
rank length \(H\). In particular \(M=H+\delta\), \(0\le\delta<1\).

Let its lower endpoint have sum \(b\). Any positive difference
of ranked sums within \([b,b+M]\) is its rank difference plus
an excess in \([0,\delta]\), so its integer part is the rank
difference. Put \(a_i=\lfloor x_i\rfloor\) and
\(\theta_i=x_i-a_i\). If the lower endpoint omits \(i<n\), add
\(i\) to that endpoint; if it contains \(i\), remove \(i\) from
the upper endpoint. This exhibits a gap of \(x_i\) within the
slab in either case. Thus \(\theta_i\le\delta\) for every \(i\),
including \(i=n\). This simple endpoint argument does not require
the projection to occupy more than half of a cube.

Delete coordinate \(n\). The projections of the slab vertices
are precisely the subsets with scores in \([b-M,b+M]\).
This band is connected: toward any target, while both types of
differing coordinate remain, add below or at \(b\) and remove
above \(b\). No step leaves the band since every weight is at
most \(M\). After either type is exhausted the path to the
target is monotone. The only point with two lifts is the common
projection of the edge endpoints, linked by the original edge.

A projected edge in direction \(i<n\) can lift to a positive
difference \(x_i\) or \(M-x_i\) in the real slab, not
\(M+x_i\). Its integer part is respectively \(a_i\) or
\(H-a_i\), since \(\theta_i\le\delta\). This is exactly the
rounded score difference and the rank difference. At the two
lifts of the repeated projection the same identity holds with
difference \(H\). It follows by connectivity that rank minus
rounded score is constant on this whole interval. Complementing
subsets gives the same property on the reflected interval;
overlap identifies the constants. Rounded scores thus increase
strictly on the entire central interval \([k,N+1-k]\).

\section{Moving a rounded collision into the central interval}
If two distinct subsets have the same rounded sum, cancel their
intersection to get disjoint sets \(U,V\). Both are nonempty since
the rounded weights are positive. Let \(C\) be their unused set
of coordinates. In real weights, take a prefix \(C'\) of \(C\)
whose sum is nearest to half the real total of \(C\).
The midpoint of the two translated scores for \(U\cup C'\) and
\(V\cup C'\) is within \(M/2\) of the real center, since
successive unused-prefix scores differ by at most \(M\).
When \(C\) is empty the midpoint is exactly the center.

The difference of the two real scores is their signed difference
of rounding errors, of magnitude at most
\(\max(|U|,|V|)\delta\). If \(C\ne\varnothing\), this is
at most \((n-2)\delta\). Each translated real score is therefore
strictly within \((H+n-1)/2\) of the center, since \(M=H+\delta\)
and \(\delta<1\). If \(C=\varnothing\), the midpoint is exact
and both scores are strictly within \((n-1)/2\) of it.
For \(n=3\) only this latter case is possible: an unused
coordinate would make both canceled sets singletons, which have
distinct floors by unit separation. In all cases the scores are
strictly within the central interval's half-width, which is at
least \(L/2\) by complement symmetry and unit spacing.
The translated sets are still distinct, and their rounded
scores are still equal. This contradicts strict increase of
rounded scores in the central interval. This proves Theorem~\ref{thm:round}. The equivalences
also use the bandwidth lower bound of Costa--Della Fiore
\cite{cdf}, noted for reals by Notik, and the inclusion of
integer configurations among real ones.

\paragraph{Consequences and verification scope.}
For \(n=9\), \((H,k,L)=(148,158,197)\); the established
integer optimum \(161\) \cite{bm} gives the earlier real
lower bound \(149\). For \(n=10\) these parameters are
\((274,288,449)\). Dyson's reported exhaustive integer search
gives \(f_{\mathbb Z}(10)=309\) \cite{dyson}, so the real lower
bound is \(275\) using that computation. We do not re-run or
formally certify Dyson's search here. For \(n=11\), the parameters
are \((526,674,701)\), giving \(f_{\mathbb R}(11)\ge527\) if
the integer maximum \(526\) is excluded. Either all-dimensional
strict bound is equivalent to the other in every dimension;
Notik asks for the real one \cite[\S7, Problem 4]{notik}.
The present theorem is a written proof, not a Lean formalization
of the uniform claim. This research note was developed with AI
assistance.

\begin{thebibliography}{9}
\bibitem{harper} L. H. Harper, Optimal numberings and isoperimetric
problems on graphs, J. Combin. Theory \textbf{1} (1966), 385--393.
\bibitem{bm} P. Borwein and M. J. Mossinghoff, Newman polynomials
with prescribed vanishing and integer sets with distinct subset
sums, \href{https://www.cecm.sfu.ca/~pborwein/PAPERS/P170.pdf}{Math. Comp. \textbf{72} (2003), 787--800}.
\bibitem{dyson} P. Dyson, Distinct subset sum search: Computation
of the 10th term of OEIS A276661 (2025).
\href{https://github.com/pwdyson/erdos_1}{Code and reported runs}.
\bibitem{cdf} S. Costa and S. Della Fiore, A bandwidth refinement
of the Erd\H{o}s distinct subset sums bound,
\href{https://arxiv.org/abs/2609.27941}{arXiv:2609.27941} (2026).
\bibitem{notik} E. Notik, Positive dual certificates for real
separated subset sums,
\href{https://arxiv.org/abs/2611.17360}{arXiv:2611.17360} (2026).
\end{thebibliography}
\end{document}
